\documentclass[11pt]{article}
\usepackage[margin=25mm]{geometry}
\usepackage{fontspec,unicode-math}
\setmainfont{Latin Modern Roman}
\setmathfont{Latin Modern Math}
\usepackage{amsmath,amssymb,mathtools}
\usepackage{xurl,hyperref}
\hypersetup{colorlinks=true,urlcolor=blue,linkcolor=blue}
\setcounter{secnumdepth}{0}
\setlength{\parindent}{0pt}
\setlength{\parskip}{5pt}
\setlength{\emergencystretch}{3em}
\newcommand{\sref}[2]{\hyperlink{source:#1:#2}{[S#2]}}
\newcommand{\cataloguescope}[1]{\paragraph{Source question.}#1}
\newcommand{\floor}[1]{\left\lfloor#1\right\rfloor}
\begin{document}
\section{Greedy Decomposition of the Spanning-Tree Polytope}\label{statement}

\subsection{Definitions and mathematical statement}

\paragraph{Shared notation.}Standard number systems are denoted by $\mathbb N,\mathbb Z,\mathbb Q,\mathbb R,\mathbb C$; index ranges, computational models and specialized notation follow the cited problem source and the definitions below.


$kP=\{kx:x\in P\}$ is a polytope dilation. An integral point of the spanning-tree polytope is the incidence vector of a spanning tree; the displayed maximum chooses the largest removable multiple of such a point while staying in the residual dilation.

A spanning tree of a connected finite graph is an acyclic connected subgraph containing every vertex. Its cost is the sum of its edge costs.

A polyhedron is an intersection of finitely many closed affine half-spaces; a polytope is a bounded polyhedron. Its faces are its intersections with supporting hyperplanes.
\paragraph{Statement recorded in the supplied source.}
Does the spanning-tree polytope $P$ admit greedy decomposition, meaning that for every positive integer $k$ and integral $w\in kP$, there is an integral $x\in P$ for which $\max\{\lambda:w-\lambda x\in(k-\lambda)P\}$ is an integer? \sref{statement}{2}
\paragraph{Source status.}
The supplied catalogue labels this record open. Its source-status metadata was checked against the supplied snapshot on 2026-10-11. This import does not independently establish current open status or mathematical correctness. The catalogue literature assessment is dated 2026-08-21.
\subsection{Short English statement}
Find a spanning tree whose maximal greedy removal from every integral dilated-polytope point has an integer coefficient.
\subsection{Sources}
\begin{itemize}
\item[\hypertarget{source:statement:1}{S1}] Ulam AI and the UnsolvedMath Contributors, UnsolvedMath: A Curated Collection of Open Mathematics Problems, catalogue entry. Catalogue curation and metadata: CC BY 4.0. The supplied source question is retained; definitions and the English summary are editorial additions.
\url{https://huggingface.co/datasets/ulamai/UnsolvedMath}.
\item[\hypertarget{source:statement:2}{S2}] Combinatorial Optimization (2012), Oberwolfach Reports, DOI 10.4171/owr/2011/53, pp. 3003--3076
\url{https://doi.org/10.4171/owr/2011/53}.
\textit{Underlying Oberwolfach Reports material retains its source-specific licence, including CC BY-SA 4.0 where applicable; see THIRD\_PARTY\_NOTICES.md.}
\end{itemize}
\paragraph{Preparation and review.}
Original question attributed to the cited source and UnsolvedMath contributors. Codex (GPT-6; exact serving revision not exposed) assisted with definitions, English summaries, source inspection and formatting. No new human mathematical authorship or independent proof verification is asserted. Catalogue identity checked against the complete supplied registry and existing repository statements; known flagged repetitions and obvious identical conjectures excluded. Full mathematical-equivalence/literature review has not been performed. Statement preparation and final review followed contribution prompts C and D; automated checks verify structure and identity, not open status.
\label{end:statement}
\end{document}
